\Needspace{16\baselineskip}
\fancyhead[L]{\small GAUSS TO VON LINDENAU.}
\begin{center}
GAUSS TO VON LINDENAU.
\end{center}

\hrule
\begin{center}
\fontsize{10}{12}\selectfont Monatliche Correspondenz. Vol. XVIII. Pp. 269..273. September 1808.
\end{center}
\hrule

\begin{flushright}
Göttingen, 30 August 1808.
\end{flushright}

\par\indent Thank you most warmly for kindly sending me the news of the two most recent comets and your observations of the new planets. It is a pity the news did not become known sooner. Of course, it is now too late to think of looking for the comets. Please allow me to share the latest results of our Juno observations in the enclosed sheet of our gel. Anz. The observations were as follows:


\begin{center}\fontsize{10}{12}\selectfont\setlength{\tabcolsep}{3pt}\renewcommand{\arraystretch}{1.15}\begin{tabular}{lrrr}
1808 & \shortstack{Mean time\\at Göttingen} & \shortstack{Apparent right\\ascension of Juno} & \shortstack{South declination\\of Juno}\\\hline
June 20 & $11^{\text{u}}\;49^{\text{m}}\;0^{\text{s}}$ & $315^\circ\;29'\;34''$ & $2^\circ\;16'\;23''$\\
22 & $12\;0\;45$ & $315\;23\;1$ & $2\;14\;29$\\
July 6 & $12\;42\;23$ & $314\;1\;0$ & $2\;19\;7$\\
30 & $10\;51\;17$ & $309\;38\;0$ & $3\;57\;13$\\
31 & $11\;40\;57$ & $309\;24\;53$ & $4\;3\;48$\\
Aug. 4 & $10\;42\;53$ & $308\;33\;44$ & $4\;30\;32$\\
5 & $10\;53\;46$ & $308\;20\;39$ & $4\;37\;36$\\
\end{tabular}\end{center}


\Needspace{11\baselineskip}
\par\indent The latest observations thus show that the error in the ephemeris has grown to $13'$ in right ascension. The error in declination has hardly changed and is still about $1'5$. No observations from elsewhere have become known so far. The last four observations give the following result for the opposition:

\begin{center}
1808 August 2. $9^{\text{u}}\,30^{\text{m}}\,43^{\text{s}}$ mean time at Göttingen\\
True longitude .\;.\;.\;.\;.\;.\;.\;.\;.\; $310^\circ\,16'\,31''5$\\
True geocentric latitude .\;.\;.\;.\; $13\;53\;56,5$ north \end{center}

\par\indent This opposition, the fourth observed so far, was then combined with those of 1804, 1806, and 1807 to determine the following new elements (VIII), which still fit all observations made to date very well, [they are reprinted above on page 306 from the Gött. gel. Anzeigen, 1808 Aug. 25.]

\par\indent We have already observed Vesta on the meridian eight times since 12 August, but some of the observations have yet to be calculated.

\par\indent [I reproduced the rest of the letter, which announces the paper on determining the altitude of the pole, above on page 146 of this volume, after that paper. \textsc{Schering}.]
